\section{Introduction}
\label{sec:intro}

Public procurement accounts for 12--15\% of GDP in OECD countries and a
substantially larger share in developing economies
\citep{oecd2019government}. Bid rigging---coordinated non-competitive
bidding among ostensible competitors---inflates procurement costs by an
estimated 10--50\%, with a median overcharge around 20\%
\citep{connor2007price}. Existing detection approaches rely either on
\emph{reactive} methods (leniency programs, whistleblowers) that depend
on cartel insiders coming forward, or on \emph{proactive} statistical
screens (variance tests, bid-rotation indices, coefficient-of-variation
thresholds) that typically require granular bid-level data
\citep{harrington2008detecting, imhof2019detecting}. In many developing
countries, however, bid-level data are unavailable or unreliable, leaving
procurement agencies with limited screening tools.

We propose a firm-level screening marker based on
\emph{frequent losers} (FL): firms that never win a single tender yet
participate in an abnormally large number of procurement processes. Under
the cover bidding hypothesis, these firms serve as ``shill bidders''
deployed by cartels to inflate the apparent number of competitors, create
the illusion of competition, and satisfy minimum-bidder requirements. The
FL screen has two practical advantages over bid-level screens: (i) it
requires only participation and outcome data, not bid values; and (ii)
it produces firm-level flags that implicate all tenders in which a
flagged firm participates, enabling targeted investigation.

This paper advances two hypotheses of increasing strength:

\begin{description}
  \item[H1 (Screening marker):] FL presence is a useful
    \emph{screening marker} for potential collusion---it correlates
    with price anomalies and overlap with known cartel participants,
    regardless of whether every FL firm is itself a cover bidder.
  \item[H2 (Causal mechanism):] FL firms are \emph{themselves} cover
    bidders deployed by cartels to inflate bidder counts---implying a
    direct causal link between FL presence and price inflation.
\end{description}

\noindent Our primary contribution is to show that FL presence is a
robust and implementable screening marker for bid rigging (H1). We then
show that a range of independent empirical tests---instrumental
variables, co-bidding network analysis, and bid-level
diagnostics---converge in a manner most consistent with a cover-bidding
mechanism (H2). We do not claim to identify all cartels in the sample,
estimate structural overcharges, or measure aggregate collusion welfare
losses. The FL screen complements existing tender-level forensic tools
\citep{imhof2019detecting} with a \emph{firm-level} criterion that is
operable with minimal data---participation records and outcome
flags---making it particularly relevant for competition authorities in
developing countries where bid-level data are unavailable.

\subsection{Contributions}

Using 4.5 million tender-items and 40 million bids from S\~ao Paulo's
Bolsa Eletr\^onica de Compras (BEC) platform (2009--2019), we make three
contributions.

\emph{First}, we show that FL presence is a reliable empirical marker
for procurement anomalies (H1). FL-present tenders exhibit 4--9\%
higher negotiated prices in OLS specifications with item, year, and
purchasing-unit fixed effects. After absorbing item and time fixed % IJIO-R1: minor-7 — Residualized gap added to Introduction
effects, the FL price gap narrows to 0.050 log points (5.0\%),
confirming that the unconditional estimate is not driven by systematic
differences in the types of goods or services procured in FL-present
tenders. External validation against CADE
(Brazil's competition authority) convictions shows that 7.1\% of FL
firms co-participate with convicted cartelists---3.5 times the rate
expected by chance after controlling for participation intensity
($p < 0.001$, permutation test). The FL--price association persists
after excluding all CADE-involved markets, indicating that the screen
detects anomalies beyond known cartels.

\emph{Second}, we provide evidence consistent with the causal mechanism
(H2) using a leave-one-out instrumental variable (IV) strategy. The
instrument exploits supply-side variation in FL firms' availability: for
each purchasing unit, we measure FL activity at \emph{other} purchasing
units in the same product market and year. The first stage is strong
($F = 396$) and the IV estimate of a 21\% price markup exceeds the OLS
estimate, consistent with measurement-error attenuation. Co-bidding
network analysis reveals that the FL--price effect is concentrated in
competitive markets (low winner concentration), where cover bidding is
most profitable, rather than in concentrated markets where a dominant
firm already controls pricing.

\emph{Third}, we provide statistical evidence of bid coordination using
\cites{bajari2003deciding} framework. FL bid residuals violate both
exchangeability (Kolmogorov--Smirnov $D = 0.15$, $p < 0.001$) and
conditional independence (pairwise product = 5.16 for FL vs.\ 2.21 for
non-FL; bootstrap 95\% CI for the difference: $[2.57, 5.21]$).

The remainder of this paper is organized as follows.
Section~\ref{sec:literature} reviews related literature.
Section~\ref{sec:conceptual} presents the conceptual framework with
testable predictions. Section~\ref{sec:data} describes the data and FL
definition. Section~\ref{sec:cade} provides external validation against
CADE convictions. Section~\ref{sec:empirical} details the empirical
strategy. Sections~\ref{sec:results} and~\ref{sec:mechanisms} present
results and mechanism tests. Section~\ref{sec:robustness} reports
robustness checks. Section~\ref{sec:conclusion} concludes.
